% EMATM0067 — Group Report
% Task 3: Comparative Corpus Analysis
% NeurIPS 2026 template

\documentclass{article}

\usepackage[final]{neurips_2026}
\usepackage[utf8]{inputenc}
\usepackage{amsmath}
\usepackage{amssymb}
\usepackage{graphicx}
\usepackage{booktabs}
\usepackage{float}
\usepackage{hyperref}

% Fix for NeurIPS 2026 undefined \@trackname in final mode
\makeatletter
\providecommand{\@trackname}{}
\makeatother

\title{Comparative Corpus Analysis of Scientific Abstracts\\
       in AI, Machine Learning and NLP}

\author{
  Ayush Bhadauria \\
  Student ID: 2685201 \\
  \texttt{zj25233@bristol.ac.uk}
  \And
  Jingting Gu \\
  Student ID: 2793772 \\
  \texttt{yn25187@bristol.ac.uk}
  \And
  Pavithraa Gangadhar \\
  Student ID: 2757887 \\
  \texttt{nc25293@bristol.ac.uk}
  \AND
  Chenyang Yan \\
  Student ID: 2765299 \\
  \texttt{th25186@bristol.ac.uk}
  \And
  Balarammahanti Siva Ganesh \\
  Student ID: 2799846 \\  \texttt{en25786@bristol.ac.uk}
  \AND
  University of Bristol \quad EMATM0067 --- Spring 2026
}

\begin{document}

\maketitle

% ============================================================
% 1. ABSTRACT
% ============================================================
\begin{abstract}
Scientific abstracts in AI, machine learning, and NLP provide
a compact record of how research priorities shift over time.
We analyse a corpus of arXiv abstracts spanning 1991--2021
across three subfields (\texttt{cs.AI}, \texttt{cs.LG},
\texttt{cs.CL}) to understand how research focus and framing
have evolved. We investigate two analytic axes: text
representation (TF-IDF unigrams, n-gram TF-IDF, and BERT
embeddings) and topic discovery method (keyword analysis,
LDA, and BERTopic). Our findings show that TF-IDF and n-gram
keyword analysis most reliably capture paradigm shifts, tracing
a clear transition from symbolic AI in the 1990s to
transformer-dominated research by 2015--2021. BERTopic reveals
that the modern era is defined by thematic diversification
rather than concentration around a single topic. LDA surfaces
coherent latent themes but is sensitive to vocabulary drift
across eras. BERT embeddings organise papers by topic rather
than era, making them less suited for temporal drift analysis.
Truncated SVD trajectory analysis confirms gradual but
accelerating vocabulary change, with the sharpest shift
occurring between 2007--2014 and 2015--2021.
\end{abstract}


% ============================================================
% 2. INTRODUCTION
% ============================================================
\section{Introduction}

The volume of scientific literature in Artificial Intelligence
(AI), Machine Learning (ML), and Natural Language Processing
(NLP) has grown at an extraordinary rate over the past three
decades. As the field has matured, the problems researchers
pursue, the vocabulary they use, and the way they frame their
contributions have changed fundamentally. Tracking these
changes at scale is difficult to do manually, but text
analytics provides a principled means of surfacing trends
across thousands of papers \citep{arxiv2020}.

This project uses the arXiv dataset available via HuggingFace,
which contains scientific abstracts annotated with metadata
including title, author, and category. We focus on three
subject categories: \texttt{cs.AI} (Artificial Intelligence),
\texttt{cs.LG} (Machine Learning), and \texttt{cs.CL}
(Computation and Language / NLP). The dataset was filtered to
these three categories to reduce noise from unrelated fields.
Publication year was extracted from each paper's arXiv
identifier, and abstracts were grouped into four periods:
1991--1998, 1999--2006, 2007--2014, and 2015--2021. These
boundaries broadly align with known transitions in the field:
the symbolic AI era, the statistical learning era, the rise
of deep learning, and the transformer era respectively.

To illustrate the vocabulary shifts present in the corpus,
consider this abstract from the earliest period:
\textit{``[PASTE ONE SENTENCE FROM A 1991--1998 ABSTRACT HERE
-- e.g. something about expert systems, logic, or knowledge
representation].''}
Contrast this with a recent abstract from 2015--2021:
\textit{``[PASTE ONE SENTENCE FROM A 2015--2021 ABSTRACT HERE
-- e.g. something about pre-trained transformers, attention,
or deep neural networks].''}
Even at a glance, the differences in vocabulary and framing
are substantial, motivating a systematic analysis.

A good solution to this task should satisfy three criteria.
First, findings should be \textit{consistent} — trends seen
in one method should appear across others, not be an artefact
of a single analytical choice. Second, outputs should be
\textit{interpretable} to a reader with domain knowledge,
with keywords and topics corresponding to recognisable
research concepts. Third, findings should
\textit{align with known patterns}, such as the emergence of
deep learning after 2012 and transformer-based methods after
2017 \citep{devlin-etal-2019-bert}.

To understand how different analytical choices affect these
criteria, we explore two experimental axes. The first is
\textbf{text representation}: we compare TF-IDF unigrams,
n-gram TF-IDF (bigrams and trigrams), and contextual BERT
embeddings. The second is \textbf{topic discovery method}:
we compare TF-IDF keyword analysis, Latent Dirichlet
Allocation (LDA), and BERTopic. We additionally apply
Truncated SVD to trace the vocabulary trajectory of the
field and t-SNE to visualise how BERT embeddings distribute
across periods. Section~\ref{sec:methods} describes the
pipeline. Section~\ref{sec:evaluation} presents results and
evaluation. We conclude in Section~\ref{sec:conclusions}.


% ============================================================
% 3. METHODS
% ============================================================
\section{Methods}
\label{sec:methods}

\subsection{Data Preparation}

The corpus was loaded from the arXiv HuggingFace dataset
\citep{arxiv2020} and filtered to \texttt{cs.AI},
\texttt{cs.LG}, and \texttt{cs.CL}. Unnecessary columns
were removed, abstracts with missing text were dropped, and
category labels were standardised. Papers were then assigned
to one of four time periods: 1991--1998, 1999--2006,
2007--2014, and 2015--2021. A pronounced class imbalance
exists across periods: Period 4 contains far more papers
than earlier periods, reflecting the exponential growth of
the field after 2012. To ensure fair comparisons, all
keyword and topic scores were normalised rather than using
raw counts.

\subsection{Text Preprocessing}

All abstracts went through the following pipeline:
lowercasing, removal of special characters and digits,
tokenisation, stopword removal using the combined
scikit-learn and NLTK English stopword lists
\citep{scikit-learn, nltk}, and lemmatisation to reduce
tokens to their root forms. A custom domain stopword list
was additionally applied to suppress high-frequency but
uninformative academic terms such as ``paper'',
``proposed'', and ``results'', which would otherwise
dominate TF-IDF scores without providing
era-discriminative signal.

\subsection{Analytic Axis 1: Text Representation}
\label{subsec:representation}

We compare three text representations, each offering a
different trade-off between interpretability and semantic
richness.

\paragraph{TF-IDF Unigrams (Baseline).}
TF-IDF weights each word by its frequency in a document
relative to the corpus, producing a sparse vector where
each dimension maps directly to a human-readable word.
\textit{Hypothesis: Unigram TF-IDF will identify
distinctive keywords per period but will miss multi-word
methodological phrases and semantically related terms.}

\paragraph{N-gram TF-IDF.}
We extend TF-IDF to bigrams and trigrams (range 2--3)
using \texttt{CountVectorizer} \citep{scikit-learn},
capturing compound phrases like ``convolutional neural
network'' that unigram analysis splits. N-gram density
is also tracked as a proxy for the technical specificity
of writing over time.
\textit{Hypothesis: N-gram TF-IDF will reveal
method-specific vocabulary shifts more clearly than
unigrams, at the cost of a larger, sparser feature space.}

\paragraph{BERT Embeddings.}
We use \texttt{bert-base-uncased} \citep{devlin-etal-2019-bert}
to produce dense 768-dimensional embeddings per abstract,
computed as the mean-pooled final hidden layer over all
tokens.
\textit{Hypothesis: BERT will capture semantic similarity
effectively but may be less sensitive to the
vocabulary-level paradigm shifts that characterise era
transitions.}

\subsection{Analytic Axis 2: Topic Discovery Method}
\label{subsec:comparison}

\paragraph{TF-IDF Keyword Analysis (Baseline).}
For each period and category we extract the top-$k$ terms
by normalised TF-IDF score at both unigram and n-gram
levels, to identify concepts gaining or losing prominence.
\textit{Hypothesis: Keyword analysis will be the most
interpretable approach but may miss latent themes spanning
multiple vocabulary items.}

\paragraph{Latent Dirichlet Allocation (LDA).}
LDA \citep{blei2003latent} models each document as a mixture
of latent topics and each topic as a word distribution.
We apply it to the preprocessed corpus to discover recurring
thematic clusters without requiring labelled data, then
track topic prevalence across periods.
\textit{Hypothesis: LDA will complement keyword analysis
by grouping co-occurring terms into coherent themes, though
its bag-of-words assumption may limit sensitivity to rapidly
evolving technical vocabulary.}

\paragraph{BERTopic.}
BERTopic \citep{grootendorst2022bertopic} clusters documents
by their BERT embeddings and labels each cluster using
class-based TF-IDF (c-TF-IDF). Topic intensity — the
relative frequency of a topic within a period and category
— is tracked over time.
\textit{Hypothesis: BERTopic will produce more semantically
coherent topics than LDA, particularly for concepts that
share meaning but vary in surface form across eras.}

\subsection{Temporal Drift Visualisation}

\textbf{Truncated SVD} \citep{scikit-learn} reduces the
TF-IDF matrix to 2D; period centroids are connected to
trace the field's macroscopic vocabulary trajectory, with
axis labels derived from the most extreme terms at each
pole. \textbf{t-SNE} \citep{vandermaaten2008} is applied
to BERT embeddings on a random sample of 5,000 abstracts
(perplexity 45) to assess whether papers cluster by time
period in semantic space.


% ============================================================
% 4. EXPERIMENTAL EVALUATION
% ============================================================
\section{Experimental Evaluation}
\label{sec:evaluation}

\subsection{Metrics and Setup}

We evaluate each approach against three criteria:
\textbf{consistency} (do findings replicate across
methods?), \textbf{interpretability} (are outputs
understandable to a domain reader?), and
\textbf{alignment} (do findings match known milestones
such as the rise of deep learning post-2012?).
Interpretability and alignment are assessed qualitatively;
we mitigate subjectivity by cross-referencing findings
across all methods. Table~\ref{tab:summary} summarises
the relative strengths of each approach.

All preprocessing artefacts were saved and reused across
experiments to guarantee consistent inputs. TF-IDF and
n-gram features used \texttt{scikit-learn}
\citep{scikit-learn}. BERT embeddings used
\texttt{bert-base-uncased} via HuggingFace
\citep{devlin-etal-2019-bert}. BERTopic used the
\texttt{bertopic} library \citep{grootendorst2022bertopic}.

\begin{table}[H]
\centering
\caption{Qualitative comparison of all methods against
evaluation criteria. \checkmark = strong,
$\sim$ = moderate, $\times$ = weak.}
\label{tab:summary}
\begin{tabular}{lccc}
\toprule
\textbf{Method} & \textbf{Consistency} &
\textbf{Interpretability} & \textbf{Alignment} \\
\midrule
TF-IDF Unigrams     & \checkmark & \checkmark & \checkmark \\
N-gram TF-IDF       & \checkmark & \checkmark & \checkmark \\
BERT Embeddings     & $\times$   & $\times$   & $\times$   \\
LDA                 & $\sim$     & $\sim$     & $\sim$     \\
BERTopic            & \checkmark & $\sim$     & \checkmark \\
SVD Trajectory      & \checkmark & $\sim$     & \checkmark \\
\bottomrule
\end{tabular}
\end{table}

\subsection{Corpus Growth and Data Imbalance}

\begin{figure}[H]
  \centering
  \includegraphics[width=0.92\linewidth]{figures/01_category_growth.png}
  \caption{ArXiv AI subfield publication growth (1991--2021).
           cs.LG dominates from 2015 onward, reaching
           approximately 37,000 papers annually by 2021.}
  \label{fig:growth}
\end{figure}

Figure~\ref{fig:growth} shows near-flat publication output
until around 2010, followed by exponential growth driven
primarily by cs.LG. This confirms the severe class imbalance
across periods: Period 4 dwarfs earlier eras by orders of
magnitude. Since raw counts would be dominated entirely by
recent papers, all keyword and topic scores are normalised
throughout the analysis.

\subsection{Axis 1 Results: N-gram Key Terms per Period}

\begin{figure}[H]
  \centering
  \includegraphics[width=0.95\linewidth]{figures/09_multi_phrases.png}
  \caption{Top 10 bigram and trigram key terms per time
           period from n-gram TF-IDF analysis.}
  \label{fig:phrases}
\end{figure}

Figure~\ref{fig:phrases} provides the clearest view of how
research language evolved. In 1991--1998, dominant phrases
include \textit{machine translation}, \textit{knowledge base},
\textit{decision tree}, and \textit{speech recognition} —
the vocabulary of symbolic AI and early NLP. By 1999--2006,
\textit{logic program} becomes the single dominant phrase,
suggesting a blend of symbolic and probabilistic approaches.
In 2007--2014, \textit{machine learn} takes over alongside
\textit{support vector} and \textit{reinforcement learn},
marking the turn to statistical methods. By 2015--2021,
\textit{machine learn} is overwhelmingly frequent
(count $>15{,}000$), joined by \textit{deep neural},
\textit{convolutional neural}, \textit{pre train}, and
\textit{reinforcement learn} — the language of the deep
learning era. This confirms our hypothesis: n-grams surface
compound methodological terms invisible to unigram analysis.

\subsection{Axis 1 Results: Keyword Evolution Heatmap}

\begin{figure}[H]
  \centering
  \includegraphics[width=0.80\linewidth]{figures/08_keywords_period_evolution.png}
  \caption{Rise and decay of AI concepts across periods
           (normalised TF-IDF). Red = high; purple = absent.}
  \label{fig:evolution}
\end{figure}

Figure~\ref{fig:evolution} tracks 30 terms across all
periods. Symbolic AI terms — \textit{grammar}, \textit{parse},
\textit{logic} — score modestly early and disappear by
Period 4. \textit{Deep} scores 0.0 in the first two periods,
rises to 52.4 in 2007--2014, and reaches 2,114.4 by
2015--2021. \textit{Network} grows from 7.5 to 3,635.9
and \textit{learn} from 19.6 to 4,080.8. The transition
accelerates sharply after 2014, consistent with the
n-gram results. Agreement between both TF-IDF methods
strengthens confidence that these are genuine field-level
shifts rather than noise.

\subsection{Axis 2 Results: LDA Topic Analysis}

We applied LDA with $k=5$ topics to the preprocessed corpus
using a bigram-aware \texttt{CountVectorizer} with a vocabulary
of 1,000 features. Table~\ref{tab:lda} shows the top keywords
for each discovered topic and our interpretation of each.

\begin{table}[H]
\centering
\small
\caption{LDA topic definitions (top keywords) and
         interpretations.}
\label{tab:lda}
\begin{tabular}{cll}
\toprule
\textbf{Topic} & \textbf{Top Keywords} & \textbf{Interpretation} \\
\midrule
1 & learn, machine, machine learn, & General ML \\
  & optimization, technique, cluster & Methods \\
2 & function, distribution, sample, & Statistical / \\
  & loss, bound, gradient, estimate & Learning Theory \\
3 & agent, policy, human, environment, & Reinforcement \\
  & user, question, reinforcement & Learning / Agents \\
4 & graph, label, representation, train, & NLP / Knowledge \\
  & domain, knowledge, speech, sentence & Representation \\
5 & network, neural, neural network, & Deep Learning / \\
  & image, deep, architecture, layer & Neural Networks \\
\bottomrule
\end{tabular}
\end{table}

\begin{figure}[H]
  \centering
  \includegraphics[width=0.95\linewidth]{figures/lda_category_period_evolution.png}
  \caption{LDA topic intensity (mean distribution) across
           categories and periods. Topic 4 (orange) grows
           steadily in cs.CL; Topic 5 (green) rises across
           all subfields from 2007 onward.}
  \label{fig:lda}
\end{figure}

Figure~\ref{fig:lda} and the underlying data reveal several
clear patterns. In \texttt{cs.CL}, Topic 4 (NLP and
knowledge representation) is the dominant topic throughout
the entire corpus and grows consistently across all four
periods, rising from an intensity of 0.41 in 1991--1998 to
0.58 by 2015--2021. This is the strongest single trend in
the LDA results, reflecting the sustained and growing focus
of the NLP community on representation learning, language
understanding, and knowledge-grounded tasks. Topic 1
(general ML methods) correspondingly declines in cs.CL
from 0.27 to 0.11 over the same period, as the subfield
became more specialised.

In \texttt{cs.AI}, Topic 3 (reinforcement learning and
agent-based systems) remains consistently prominent
throughout all periods, peaking at 0.39 in 1999--2006 and
staying near 0.36 by 2015--2021. This reflects the
longstanding centrality of agent-based reasoning to the AI
subfield, even as neural methods rose elsewhere. Topic 5
(deep learning and neural networks) grows from just 0.04 in
1991--1998 to 0.15 by 2015--2021, consistent with the
broader field-wide trend.

In \texttt{cs.LG}, Topic 2 (statistical and mathematical
learning theory) shows a notable spike, rising sharply from
0.11 in 1991--1998 to 0.31 by 1999--2006 and remaining
elevated through 2007--2014. This aligns with the statistical
learning era and the growth of kernel methods and
theoretical ML research in that period. By 2015--2021,
Topic 5 (neural networks) grows to 0.21, reflecting
the shift of cs.LG toward deep learning.

These findings partially confirm our hypothesis. LDA
successfully identified coherent, interpretable thematic
groupings that align with known developments in each
subfield. However, the dominance of the term ``learn''
across all topics — a consequence of it surviving the
stopword filter — reduces the discriminative sharpness of
topic boundaries. This is a known limitation of LDA applied
to domain-homogeneous corpora: when a field is uniformly
defined by a core concept such as learning, that concept
saturates the topic-word distributions and limits the
separation between topics compared to what BERTopic
achieves through embedding-based clustering.

\subsection{Axis 2 Results: BERTopic Evolution}

\begin{figure}[H]
  \centering
  \includegraphics[width=0.95\linewidth]{figures/bert_category_period_evolution.png}
  \caption{BERTopic topic intensity across categories and
           periods. Topic 1 (teal) peaks in 2007--2014 and
           declines sharply by 2015--2021 in all subfields.}
  \label{fig:bert}
\end{figure}

Figure~\ref{fig:bert} shows Topic 1 dominating the early
and middle periods, peaking at 0.31 in cs.CL (2007--2014)
and 0.30 in cs.LG (1999--2006), before collapsing below
0.05 by 2015--2021. Rather than a failure of BERTopic,
this reflects genuine field-level diversification: as
publication volume expanded exponentially, research spread
across computer vision, reinforcement learning, language
modelling, and graph networks simultaneously. This
finding complements the n-gram results and aligns with
our hypothesis.

\subsection{Temporal Drift: SVD and t-SNE}

\begin{figure}[H]
  \centering
  \includegraphics[width=0.78\linewidth]{figures/04_tfidf_trajectory.png}
  \caption{Truncated SVD projection. Period centroids trace
           a rightward trajectory along the ``buyer
           $\to$ learn'' semantic axis.}
  \label{fig:trajectory}
\end{figure}

Figure~\ref{fig:trajectory} shows period centroids moving
progressively rightward along the primary SVD axis, labelled
from ``buyer'' toward ``learn''. The trajectory is compact
overall, but the largest single step occurs between Period 3
and Period 4, consistent with the accelerating shift seen in
the keyword heatmap. The ``learn'' pole directly reflects
the first principal component of 30 years of
TF-IDF-weighted vocabulary variation.

\begin{figure}[H]
  \centering
  \includegraphics[width=0.76\linewidth]{figures/03_embeddings_clustering.png}
  \caption{t-SNE of BERT embeddings (5,000 abstracts).
           Periods are interspersed — BERT organises
           by topic, not era.}
  \label{fig:tsne}
\end{figure}

Figure~\ref{fig:tsne} shows BERT embeddings do not cluster
by time period. A plausible explanation is that BERT encodes
research content at the topic level — reinforcement learning
appears in every era — rather than reflecting the
vocabulary-level shifts TF-IDF captures. This is a useful
negative finding: for diachronic vocabulary analysis,
surface-level representations appear more sensitive to era
transitions than deep contextual ones.

\subsection{Error Analysis and Discussion}

The custom stopword list required manual judgement: terms
like ``learn'' and ``network'' carry genuine discriminative
signal but their extreme frequency in Period 4 may mask
subtler thematic patterns. This directly affected LDA —
the term ``learn'' saturates all five topic-word
distributions, reducing boundary sharpness between topics.
A stricter stopword policy or larger vocabulary
($>1{,}000$ features) would likely produce cleaner topic
separation. BERTopic's intensity collapse in 2015--2021
confirms diversification but cannot describe the content
of that diversity without examining full topic label lists.
The t-SNE visualisation is affected by class imbalance —
Period 4 dominates the 5,000-abstract sample — meaning
earlier periods are underrepresented and harder to
interpret geometrically.

Despite these limitations, findings are highly consistent
across methods. TF-IDF keyword and n-gram analysis produced
the most interpretable and robust outputs. BERTopic added
a complementary view of thematic diversification. LDA
provided interpretable topics but showed less stability
across hyperparameter choices. BERT embeddings were least
effective for era-level temporal analysis. The main
limitation of the overall approach is treating categories
as static and mutually exclusive, whereas in practice many
papers span subfields and period transitions are gradual.


% ============================================================
% 5. CONCLUSIONS
% ============================================================
\section{Conclusions}
\label{sec:conclusions}

This work applied a multi-method text analytics pipeline to
30 years of arXiv abstracts from AI, ML, and NLP subfields.
Across all methods, the findings tell a consistent story:
the corpus transitions from symbolic AI and formal methods
in the 1990s, through probabilistic and statistical learning
in the 2000s, into neural methods from 2007, and finally
into a highly diverse, transformer-dominated landscape from
2015.

TF-IDF keyword and n-gram analysis were the most effective
tools for detecting this transition. Multi-word phrases
such as \textit{convolutional neural}, \textit{pre train},
and \textit{reinforcement learn} were essential signals
invisible to unigram analysis. BERTopic added a distinct
insight: the modern era is characterised by the
fragmentation of research into many specialised directions,
reflected in sharply declining per-topic intensity after
2014. LDA surfaced interpretable thematic groupings that
broadly aligned with keyword trends, though it was more
sensitive to hyperparameter choices than BERTopic.

BERT embeddings did not cluster by time period, which is a
useful negative result. Semantic representations encode
research content at the topic level rather than the era
level, making TF-IDF-based representations more appropriate
for diachronic vocabulary analysis of this kind.

Open challenges include addressing class imbalance more
principally, developing quantitative interpretability
metrics for topic model outputs, and extending the analysis
beyond abstracts to full paper texts or citation networks,
which could reveal influence patterns invisible at the
abstract level.

Code repository: \url{https://github.com/[YOUR REPO LINK]}


% ============================================================
% REFERENCES
% ============================================================
\bibliographystyle{unsrt}
\bibliography{references}

\end{document}